\section{Introduction}
\vspace{-3pt}

\label{sec:intro}
Diffusion models have emerged as a powerful generative paradigm across diverse domains, including image synthesis~\citep{ddpm, LDM}, biological sequence generation~\citep{diff-bio, diff-protein}, and language modeling~\citep{SEDD, D3PM, MDLM}.
Unlike autoregressive~(AR) models, which generate sequences via a left-to-right factorization, diffusion models formulate generation as an iterative refinement process that transforms samples from a simple initial distribution toward the target distribution through learned backward transitions, and therefore have the potential to improve long-horizon text modeling and accelerate generation.

Conventional diffusion models typically prescribe a forward noising process that gradually perturbs data samples under a predefined noise schedule~\cite{ddpm, MD4} from the initial distribution.
This forward process defines the intermediate distribution path that the backward model learns to invert, thereby strongly shaping the difficulty of both training and inference~\citep{D3PM,iDDPM,design-path}.
In continuous domains, this paradigm is naturally supported by Gaussian perturbations, which exploit local neighborhoods under Euclidean geometry and preserve meaningful proximity between nearby states~\citep{ddpm, score-based}.
In discrete language modeling, however, there is no analogous canonical geometry for noise design.
Token indices carry no intrinsic metric, and semantic proximity is highly context-dependent~\citep{D3PM, ctx_dep}.
Consequently, the choice of initial distribution and forward corruption process becomes central to modeling language with diffusion: a poorly matched forward process may induce intermediate states that are misaligned with the drafts and errors encountered during generation, making the learned backward transitions less effective~\citep{SCUD, mdpo}.

Many structured perturbation approaches have been proposed for diffusion language modeling, including structured discrete transition kernels~\citep{multinomial_diffusion, ctdd, D3PM}, absorbing-state masking~\citep{MDLM,MD4}, discrete score-based objectives~\citep{SEDD, MDM_secret}, and scaling recipes adapted from autoregressive models~\citep{DiffuLLaMA,LLaDA, sdar, dream7b}.
However, discrete diffusion methods with predefined corruption paths still exhibit a significant performance gap relative to AR models in language modeling~\citep{D3PM,gulrajani2023likelihood,he2023diffusionbert,SEDD}, showing that without a natural neighborhood structure in discrete language spaces, the best perturbation design remains unclear.

Motivated by the difficulty of designing forward perturbations in discrete space, we ask whether diffusion language modeling can bypass a hand-designed forward corruption process altogether.
We propose \textbf{F}orward-F\textbf{Re}e \textbf{D}iffusion L\textbf{A}nguage Model~(\method) to answer this question.
We emphasize that diffusion language models recursively refine a distribution from an initial draft distribution and progressively close the gap to the target data distribution through learned reverse transitions.
Based on this view, \method uses drafts generated by the current model as implicit intermediate states, allowing the distributional path to be induced by the sampler itself rather than prescribed by an external noising schedule.
This design makes \method both {\bf neighborhood agnostic}, \ie, avoiding artificial token-level perturbations in discrete space, and {\bf sampler adaptive}, \ie, training each transition on the drafts and errors produced by the current sampler.
\method naturally supports flexible refinement parameterizations, and we instantiate it with two designs: {\bf 1)} self-refinement, which directly revises the model's own draft; and {\bf 2)} stochastic Best-of-$N$ refinement, which proposes multiple candidate sequences in parallel and uses an additional score head to select the most promising one.

The recurrent use of the same refinement backbone also places \method directly in the Looped Transformer family~\citep{dehghani2019universal,saunshi2025latentthoughts,geiping2025recurrentdepth}.
Specifically, \method defines a new draft-space form of Looped Transformer: each pass revises an explicit model-generated token draft and feeds the revised draft into the next pass, rather than carrying an internal latent state across loops.
Because every loop produces a candidate output, it can be supervised directly against the clean target.
During training, the preceding refinement passes are detached and gradients are applied only through the final pass, so \method avoids BPTT across refinement iterations while retaining ordinary backpropagation within the final Transformer pass.
This draft-space loop also preserves block-parallel token updates and naturally supports branching into parallel candidates for Best-of-$N$ refinement.
We discuss this connection further in Sec.~\ref{sec:looped_connection}.

% At inference time, \method repeatedly proposes candidate sequences and refines its own drafts, aligning training with the generation dynamics rather than predefined corruptions.
% This also enables flexible inference-time scaling, allowing early exit from intermediate marginals with limited computation or additional refinement steps to improve sample quality.

We train \method at the 4B scale using a 10B-token continued-pretraining corpus, which is substantially smaller than the training budgets used by several competing diffusion baselines.
We evaluate it extensively across general knowledge, mathematical reasoning, and code generation benchmarks in the sub-8B regime.
Although substantially smaller than several competing diffusion models, \method-4B {\bf surpasses 7--8B} diffusion base models on four reasoning and coding benchmarks, achieving absolute improvements of 5--15\%, while remaining competitive on general knowledge tasks.
\method also establishes a stronger quality--speed Pareto frontier, requiring fewer forward passes to reach comparable accuracy and generating more tokens per forward pass than existing diffusion baselines, with an average speedup of 1.5--1.8$\times$.
Additional analyses further show that \method scales effectively with refinement computation, yielding consistent gains over refinement iterations.

Our main contributions are as follows:
(1) We propose \method, a forward-free diffusion language model that uses model-generated drafts as implicit intermediate states and learns to refine them through flexible parameterizations. The recurrent use of a shared refinement backbone makes \method a draft-space Looped Transformer whose training avoids BPTT across refinement iterations.
(2) We empirically evaluate \method across general knowledge, mathematical reasoning, and code generation benchmarks, demonstrating that \method-4B substantially outperforms 7--8B diffusion baselines and establishes a strong quality--speed Pareto frontier.
